\chapter{Critical analysis of security properties}
\label{chap:sicurezza}

\section{Attack scenarios and mitigations}
\label{sec:scenari-attacco}

The previous chapters describe RAM-USB's defenses from the point of view of those who build them.
This section rereads them from the point of view of those who try to compromise them.

\subsection{Email address enumeration at registration}

\texttt{POST /api/users} is a public endpoint that responds 201 if the submitted email is available
and 409 if it is already registered, without specifying whether the conflict comes from the email or from the
SSH key. However, anyone who submits a client-validated password and a freshly generated SSH key,
and therefore one that is certainly unique, has already ruled out that the conflict is on the SSH key.

\autoref{sch:enumerazione} summarizes the procedure, with the following symbols:

\begin{itemize}
    \item $e$: the candidate email address;
    \item $p$: a client-validated password;
    \item $k$: an SSH public key generated by the attacker;
    \item $r$: the HTTP status code returned by \texttt{POST /api/users}.
\end{itemize}

\begin{schema}
\centering
$\begin{array}{rl}
1. & k \leftarrow \mathrm{KeyGen}() \\[4pt]
2. & r \leftarrow \mathrm{POST}\ \texttt{/api/users}(e, \, p, \, k) \\[4pt]
3. & r = 409 \;\Longrightarrow\; e \mbox{ already registered}
\end{array}$
\caption{Checking whether an email address is already registered.}
\label{sch:enumerazione}
\end{schema}

Login, on the other hand, hides this information much better, returning the same 401 both for a nonexistent
email and for a wrong password, with the same log and the same response time
(\autoref{sec:login}). That protection was never extended to registration.

The registration endpoint also has a side effect far more serious than enumeration, caused
by the fact that the system does not verify that whoever submitted the registration credentials actually
owns them.

A malicious actor who registers with an email that does not belong to them permanently prevents registration
by its legitimate owner, forcing the owner to use a different email address.

Both problems could be solved by responding to registration requests identically in
all cases, for example always with a 202, while the actual outcome would travel out-of-band, in a
confirmation email sent to the given address.

This, however, requires rethinking how to deliver to the user the pre-auth key that currently travels inside the
201 registration response.

\subsection{Administrative access to credentials}
\label{sec:accesso-credenziali}

Email and password pass, encrypted in transit, through Entry-Hub, Security-Switch, and Database-Vault.
Whoever controls even just one of these processes can read them at the point where they are in plaintext in memory.

At rest, however, the password is never recoverable, not even by an administrator with full
access to the database. The email, instead, is recoverable by an administrator who holds the master key.

The problem of credentials captured in transit could be mitigated by encrypting the credentials end-to-end
between the client and Database-Vault, but in this case the problem would still not be completely
solved, and additional complexity would be introduced into the system. The only real solution is to use a zero-knowledge
authentication system.

\subsection{Offline address verification from a database copy}

An attacker who holds a stolen copy of the database can immediately check
whether a suspected email address is present in the database or not. The attacker only needs to compute the SHA-256 hash
of the candidate address and compare it with the hashes stored in the database.

Extracting an arbitrary email from the database, however, requires a very different effort.
The attacker cannot decrypt the email without the master key and is therefore forced to attack its SHA-256
hash.
The attacker could, however, identify some emails if they appeared in public dictionaries, using a
dictionary attack; otherwise, the only viable path remains a brute-force attack.

\autoref{sch:bruteforce} formalizes the procedure, with the following symbols:

\begin{itemize}
    \item $D$: the set of digests present in the stolen database;
    \item $W$: a dictionary of plausible addresses;
    \item $e$: an email address drawn from $W$.
\end{itemize}

\begin{schema}
\centering
$\begin{array}{rl}
1. & \forall e \in W: \\[4pt]
2. & k_h \leftarrow \mathrm{SHA\mbox{-}256}(\mathrm{lower}(e)) \\[4pt]
3. & k_h \in D \;\Longrightarrow\; e \mbox{ already registered}
\end{array}$
\caption{Search for registered addresses from the stolen database, by dictionary or brute force.}
\label{sch:bruteforce}
\end{schema}

Both of these vulnerabilities can be mitigated by replacing the plain digest with an HMAC-SHA256 keyed with
a dedicated pepper, generated with sufficient entropy and distinct from the one used for password hashing.
An attacker who steals the database does not automatically obtain the pepper as well, and therefore cannot
recover the email addresses in a reasonable time.

The theft of user credentials, although very serious, does not compromise the backups and does not prevent
the user from accessing them. Accessing them requires the SSH private key, which the user keeps locally.
Nor are the stolen credentials enough to delete or modify the user's account, because the system does not
guarantee users the right to be forgotten.

\section{Failure isolation}

This section analyzes which of the User Requirements remain guaranteed even in the worst case: the
full compromise of a service.
It also analyzes how to resolve the compromise where possible, provided that it is detected.

\begin{table}[H]
\centering
\caption{User Requirements verified in this section against the full compromise of a service.}
\label{tab:ru-isolamento-fallimenti}
\begin{tabular}{@{}l p{\dimexpr0.75\linewidth\relax}@{}}
\hline
\textbf{Requirement} & \textbf{Description} \\
\hline
RU-04 & Be able to access one's own backup only after having authenticated. \\
RU-05 & Have the guarantee that no one, administrators included, can read one's own backups. \\
RU-06 & Have the guarantee that no one, administrators included, can read one's own
access credentials. \\
RU-08 & Have the certainty that one's own data is isolated from that of other users. \\
\hline
\end{tabular}
\end{table}

The remaining properties fall outside this verification: RU-01, RU-02, RU-03, RU-07, and RU-10 are functional
or availability requirements that a compromised service easily denies, and therefore do not require
a case-by-case verification; RU-09 remains out of scope as stated in \autoref{chap:requisiti}.

\subsection{Entry-Hub}

An attacker who gains control of Entry-Hub can read, block, or tamper with every registration and login
request. They can also abuse its reachability toward Security-Switch,
Certificate-Authority, and the MQTT broker.

Tampering, however, does not affect RU-05, because the Restic repository password originates on the user's
device and never passes through Entry-Hub, before, during, or after tampering (\autoref{chap:storage}).
Nor does it affect RU-04, because an already authenticated user reaches Storage-Service directly over the
mesh, without ever going through Entry-Hub (\autoref{chap:architettura}).
The attacker can, however, block or alter new registration and login requests.

The compromise, if detected, can easily be contained. Network-Manager can revoke
a node's membership in the mesh, cutting the compromised Entry-Hub instance off from the network, and while
the cause of the compromise is being investigated, a new Entry-Hub instance can be started in a KVM on
Proxmox, ideally on a different node (\autoref{sec:docker-vs-proxmox}). In this way, registration of new users remains impossible, since the KVM might
not have a reachable public address, but user login can take place without problems,
since it happens on the handler reachable only from the private mesh network.


\subsection{Certificate-Authority}

An attacker who gains control of Certificate-Authority could sign a valid mTLS certificate
for any \texttt{organization} (\autoref{sec:pki}) and thus impersonate any internal
service, or refuse to issue updated certificates to services that request a renewal.

Realistically, however, they cannot contact anyone other than the MQTT broker, because according to Headscale's ACL rules, the CA can only
respond to calls from internal services, without being able to initiate new communications (\autoref{sec:trust-zones}).

This does not affect RU-05 or RU-06, since the CA plays no role in backup encryption
(\autoref{sec:backup-restore}) nor in the derivation of the credential keys
(\autoref{chap:autenticazione}).
RU-08 also remains guaranteed, as it has no connection with the CA.
What breaks down is mutual authentication between services.
The services, however, being unable to renew their certificates, can no longer communicate, so
once the certificates expire, the entire system comes to a halt.

The failure of authentication between services can easily be detected.
When this happens, until the investigation into the compromise is concluded, a new instance of the
CA can be restarted in a KVM on Proxmox, ideally on
a different node, ensuring greater isolation than that offered by an LXC (\autoref{sec:docker-vs-proxmox}).


\subsection{Database-Vault}

An attacker who gains full control of Database-Vault obtains the pepper and the master key, and with
the latter decrypts every email stored in the database.
However, they do not obtain the password with the same ease. For that, the attacker has two possible paths:
the first is to steal the database and attack it offline, at the cost of Argon2id's computation time (\autoref{tab:phc}).
This approach normally requires a great deal of time and considerable computational resources, making it
practically infeasible.
Without the pepper to protect them, however, weak and very common passwords can easily be cracked with
a dictionary attack, even if protected by a slow and expensive hashing algorithm such as Argon2id.

The second path is far more insidious: the attacker can remain passive and observe the calls to Database-Vault.
In this way, not only can they read the credentials of new users, but, by analyzing login requests and
following them all the way to the database, they can identify which hashes the passwords sent by users correspond to.
This requires the malicious actor to remain hidden for the time needed to receive login
requests from all the users who registered before the infiltration.

RU-06 therefore no longer holds at all for the email, and is reduced to a time cost for complex and uncommon
passwords.

RU-05 nonetheless remains guaranteed: Database-Vault has no access to Storage-Service's filesystem; it only has the
ability to ask it to create an account (\autoref{chap:architettura}). RU-08 also remains guaranteed, for
the same reason.

In this case, even if the active attack were detected, there would be nothing left to do; one must assume
that the database has already been stolen.
If instead it were the passive attack that was detected, while the compromise is being investigated,
the registration and login services can be shut down completely, and users
who are already authenticated can be granted a reachability tag toward Storage-Service that lasts longer than the twelve-hour one they
already hold (\autoref{chap:rete}). In this way, logged-in users can continue to manage their
backups without having to re-authenticate.

In both cases, however, it would be more sensible to shut down the service and declare the incident.

\subsection{Network-Manager}

An attacker who gains control of Network-Manager can grant themselves reachability toward every other
service and toward the users (\autoref{chap:rete}). They can also cancel or extend the access grants toward
Storage-Service (\autoref{lst:grants-schema}).
They can therefore block and control all of the system's traffic, without, however, being able to read or decrypt it.
RU-04 is therefore compromised.
They still cannot access users' data or their credentials (\autoref{chap:autenticazione}), so
RU-05, RU-06, and RU-08 remain valid.

In case of compromise, while the causes are being investigated, Network-Manager can still be shut down and
a new instance started in a different KVM, ideally on a different node.

\subsection{Storage-Service}

An attacker who gains control of Storage-Service inherits the root privilege with which the
service must run (\autoref{sec:isolamento}), and with it reads every user's subdirectory
directly from the host filesystem: here RU-08 is violated.

RU-04 is violated as well, because the attacker can prevent any communication with Storage-Service.
They can even encrypt, corrupt, or delete the backups of all users, but this is a violation of RU-09,
declared out of scope for this project.
RU-05 nonetheless remains guaranteed, because the data is encrypted client-side and never arrives in plaintext on
Storage-Service (\autoref{sec:backup-restore}).
RU-06 also remains guaranteed, because access credentials do not pass through Storage-Service.

In this case, when the compromise is detected, there is nothing else to do but declare
the incident and shut down the service.
